\section{Introduction}
\label{sec:intro}

Medicinal plants are identified in the field mostly from their leaves, and errors are not equally
costly: confusing two culinary herbs is harmless, but confusing a remedy with a toxic look-alike, such
as \textit{Nerium oleander} with other Apocynaceae, can be dangerous~\citep{farkhondeh2020oleander}.
Deep networks report 97.5--99.7\% accuracy on Indian medicinal leaf
benchmarks~\citep{roopashree2021deepherb,ffncacnn2025}. Those numbers, however, are obtained on
single leaves photographed on white backgrounds, typically with random image-level splits. They say
little about how the models behave when the same species is photographed in the field, against soil,
tiles or cloth, with a phone camera, which is where such a classifier would be used.

Botanists describe a leaf by its geometry: how wide the blade is along its length, where it tapers,
how ragged its margin is, how its veins branch. These characters are measured \emph{at particular
places} on the leaf. Standard CNN and ViT classifiers average their features over the whole image and
have no built-in notion of where the leaf is, let alone of localised geometry. A natural hypothesis is
therefore that architectures which measure leaf geometry locally, and look at the image where they
measure, should transfer better from studio to field. Recent work fuses topological and morphological
summaries with deep features~\citep{phgnet2023,tdavit2026,gu2026multifiltration,blaser2026encoding},
but usually as global descriptors, and without controls that would tell whether the geometry itself
is what helps.

We test this hypothesis with a ladder of matched models that share backbone, data, vein head and
training, and differ in one ingredient at a time: (1)~a plain classifier; (2)~the same features pooled
over the leaf mask, i.e.\ knowing \emph{where} the leaf is; (3)~localised pooling over strips of the
image or over slabs of an object-centred directional sweep; and (4)~\emph{sweep-valuation attention}
(\SVA{}), which we introduce as the most structured member of the family: for each slab of the sweep it
measures how much area and boundary length a learned, appearance-weighted cubical complex of the blade
and venation adds there, and lets that measurement attend to the image patches of the same slab. A
global fusion of the same measurements tests whether geometry helps without localisation.

The answer is clear and partly negative. On ConvNeXt-T the studio-to-field gain over a plain model
is large (14 points) but splits into two parts of similar size: knowing where the leaf is, and pooling
locally within it. Measured geometry, whether fused globally or attached to localised tokens, adds
nothing detectable. The gain is small on EfficientNet-B0, absent on a ViT, and irrelevant for a
biology foundation model, which transfers far better on its own.

\paragraph{Contributions.}
\begin{enumerate}
\item \textbf{A controlled decomposition of studio-to-field transfer} for medicinal leaves on four
  backbones, with 5 folds $\times$ 5 seeds for the main backbone, pre-specified contrasts, corrected
  resampled tests with Holm correction and a bootstrap over runs and field species. Leaf location and
  localised pooling each add 5--8 points on raw field photographs of ImageNet ConvNeXt-T; the gains
  hold without rotation augmentation and add to strong background-compositing augmentation.
\item \textbf{Sweep-valuation attention}, a differentiable, backbone-agnostic module that turns a
  directional filtration of appearance-weighted cubical complexes into slab tokens, and a careful
  negative result: its measured geometry does not improve transfer beyond localised pooling, its
  slab-restricted attention can be replaced by global attention, and Euler-characteristic tokens make
  models sensitive to topology but transfer worse.
\item \textbf{A leakage-audited protocol and probes} for the Medicinal Leaf Dataset: duplicate and
  re-shot grouping, a corruption suite (MedLeaf-C), a leaf cue-conflict set, a perceptually matched
  topology-destruction test and a verified mapping of 20 DIMPSAR field species.
\end{enumerate}
Code, splits, probes, statistics and the outputs of all training runs are released.
